\documentclass[10pt]{article}
\ifdefined\TMLRSubmission
\usepackage{tmlr}
\else
\usepackage[margin=1in]{geometry}
\usepackage[numbers,square]{natbib}
\fi
\usepackage{amsmath,amssymb,amsthm,booktabs}\usepackage{lmodern}\usepackage[T1]{fontenc}
\usepackage[hidelinks,hyperfootnotes=false]{hyperref}
\ifdefined\TMLRSubmission
\hypersetup{pdftitle={Soluble State Tracking in Materialized Fixed-Point GDN+},pdfauthor={},pdfsubject={Materialized fixed-point expressivity}}
\else
\hypersetup{pdftitle={Soluble State Tracking in Materialized Fixed-Point GDN+},pdfauthor={Tarang Khanna},pdfsubject={Materialized fixed-point expressivity}}
\fi
\newtheorem{theorem}{Theorem}
\newtheorem{lemma}{Lemma}
\newtheorem{proposition}{Proposition}
\theoremstyle{remark}
\newtheorem*{remark}{Remark}
\newcommand{\Z}{\mathbb Z}
\newcommand{\D}{\mathbb D}
\newcommand{\LN}{\operatorname{LN}}
\newcommand{\relu}{\operatorname{ReLU}}
\newcommand{\Sol}{\mathsf{Sol}}
\newcommand{\Aprd}{\mathsf{Aprd}}
\newcommand{\Fin}{\mathsf{Fin}}
\newcommand{\GDNf}{\mathsf{GDN}^{+}_{f}}
\DeclareMathOperator{\R}{R}

\title{Soluble State Tracking in Materialized Fixed-Point GDN$^{+}$}
\ifdefined\TMLRSubmission
\author{Anonymous authors}
\else
\author{Tarang Khanna\\
Independent Researcher\\
\texttt{tarang.research@gmail.com}}
\fi
\date{}
\newcommand{\paperbibitem}[2]{\bibitem[#1]{#2}}

\begin{document}
\ifdefined\TMLRSubmission
  \maketitle
\else
  % Compact title for the identified preprint.
  % The official TMLR title layout is deliberately left untouched.
  \vspace*{-1.5em}
  \begin{center}
    {\Large Soluble State Tracking in Materialized Fixed-Point GDN$^{+}$\par}
    \vspace{0.7em}
    {\normalsize Tarang Khanna\par}
    {\small Independent Researcher\par}
    {\small\href{mailto:tarang.research@gmail.com}{\texttt{tarang.research@gmail.com}}\par}
  \end{center}
  \vspace{0.35em}
\fi
\begingroup\renewcommand{\thefootnote}{*}\footnotetext{Disclosure of AI assistance: large language models contributed to developing the mathematical constructions, implementing computational verification, checking mathematical arguments, and drafting this manuscript. The human author is responsible for the submitted claims and materials.}\endgroup

\begin{abstract}
Every finite soluble monoid admits a fixed-width, fixed-depth GDN$^{+}$ network
computing nonempty prefix products in the materialized fixed-point model defined
here, with constant gates $\alpha=\beta=1$, zero initial memory, and either of
two specified centered normalizations. Building on StateBench's aperiodic and
rounded-parity constructions, we supply arbitrary-modulus counters and
full-network soluble composition. Individually rounded materialized products
implement bounded counters. Shifted copies yield constant-norm phase codes that
survive normalization and residual addition; bias-free ReLU selectors return
reusable labels. A rounded zero-state marker replaces a beginning-of-sequence
token, and a parity-driven buffer supplies the exclusive prefix needed for
wreath composition. Exhaustive certificates check complete modulo-three and
$\Z_3\wr U_2$ networks under both normalizations; the general result follows
by algebraic induction. The construction depends on this evaluation program:
the tested rounded rank-one evaluator fails with unchanged weights, without
proving impossibility for rank-one models.
\end{abstract}

\section{Introduction}

State tracking asks a sequence model to apply successive updates without
losing the information needed by later ones. Finite monoid word problems
separate immutable assignment, mutable assignment, self-reference and general
finite-state updates. In StateBench, \citet{statebench} organize these as the
strict hierarchy r-trivial $\subsetneq$ aperiodic $\subsetneq$ soluble
$\subsetneq$ finite. Its Theorem~1 places fixed-precision GDN$^+$ above the
aperiodic class, and Lemma~3 gives a rounded parity orbit. The exact placement
relative to soluble tracking remains an open question in that paper.

The issue is not only whether a recurrent head can count. Its state must pass
through legal projections, normalization, residuals and feedforward layers,
then supply the timing and precision interfaces needed for composition. We
construct these interfaces in a materialized rational-grid model.

\begin{theorem}[Materialized fixed-point model]\label{thm:main}
Fix either centered normalization and the scalar operations of
\S\ref{sec:model}, and permit constant gates $\alpha=\beta=1$ as actual
parameters. For every finite soluble monoid $M$, there exist fixed finite
width, depth, head count and feedforward width, integers $D,p$, and rational
parameters in the single datatype $\D_{p,D}$ such that the zero-initialized
network returns the product of every nonempty word over $M$. Its canonical
output at every position encodes the inclusive prefix product. All scalar
parameters and pre-rounding intermediate results remain in range for arbitrary
input length. There are no biases, positional inputs or beginning-of-sequence
token. The two normalization conventions use different constructions.
\end{theorem}
The theorem establishes containment of $\Sol$ in the defined model. The proof
uses the usual Krohn--Rhodes induction; its architecture-specific content is in
three primitives and their interfaces:
\begin{enumerate}
\setlength{\itemsep}{2pt}
\item \textbf{Rescaled modular counters.} Separately rounding the $m^2$ matrix
products implements a bounded counter for every fixed $m$. Rescaling the state
and value makes the same integer transition available at compatible common
 denominators (\S\ref{sec:counter}).
\item \textbf{Normalization-exact carriers.} Shifted counter copies give a
constant-norm phase code. Opposite shared coordinates cancel in the residual;
a finite bias-free ReLU selector then returns a canonical label. A rounded
zero-state marker initializes the shifts without a BOS token
(\S\ref{sec:interface}; Figure~\ref{fig:z3}; Table~\ref{tab:example}).
\item \textbf{Composable timing.} A parity-driven buffer exposes the previous
lower prefix needed by a wreath product. Payload-preserving recoding and a
single datatype close the induction (\S\ref{sec:wreath}). Delayed composition
itself is prior work~\citep[Lemmas~18, 20]{sarrof}; the construction here realizes
it with the stated GDN$^+$ operations.
\end{enumerate}

\section{Model and arithmetic}\label{sec:model}

We specify the following operational model, following the sublayer structure
of Definitions~9, 10 and~13 in~\citet{statebench}.  Tokens are embedded and normalized, $h^0_i=\LN(E[w_i])$. Each layer computes
\[
\hat h=\LN(h),\qquad
g=h+\LN\!\big(O\cdot \mathrm{head}(\hat h)\big),\qquad
h'=g+W_2\,\relu(W_1\LN(g)),
\]
Here $\mathrm{head}$ denotes the concatenated outputs of a fixed finite
number of heads. Queries, keys and values are linear in $\hat h$, with no biases. The feedforward write is not normalized (Definition~10). A GDN head with one value column block maintains
\[
A_t=I-\beta_t k_tk_t^\top,\qquad S_t=\alpha_t A_tS_{t-1}+\beta_t k_tv_t^\top,\qquad o_t=S_t^\top q_t,
\]
with $S_0=0$. Throughout we permit constant $\alpha_t=\beta_t=1$; with unit keys the
unrounded $A_t$ is a positive-semidefinite projector. This does not assert a
spectral property of the rounded nonlinear transition map. The final output is $\arg\max h^d_nU$.

\paragraph{Arithmetic.} Numbers live in $\D_{p,D}=\{j/D:\ j\in\Z,\ |j|\le 2^p\}\cup\{\bot\}$.
This rational grid makes the evaluation order discussed in StateBench
Appendix~A.4 explicit; identification with every inherited cast convention is
not assumed. Before rounding, every scalar operation is checked against the physical
range $[-2^p/D,2^p/D]$; overflow produces absorbing $\bot$. Sums are exact
unless they overflow. Host integers used to evaluate rational expressions are
not themselves values in this datatype. A product $(a/D)(b/D)$ is rounded to $\R(ab/D)/D$ with $\R(x)=\lfloor x+\tfrac12\rfloor$ (ties up). Iterated sums are evaluated left to right. Matrix products are evaluated entrywise with each scalar product rounded, and $A_t$ is materialized before $A_tS_{t-1}$ is formed. Division and square root inside layer normalization are rounded to the grid in the same way; the construction arranges that every square-root argument is exactly $1$, so only $\sqrt{1}=1$ is used.

\paragraph{Normalization.}
For width $n$, let $\bar x=n^{-1}\sum_jx_j$. We construct different networks for
\[
\LN_{\mathrm E}(x)=\frac{x-\bar x\mathbf1}
 {\sqrt{\sum_j(x_j-\bar x)^2}},\qquad
\LN_{\mathrm V}(x)=\frac{x-\bar x\mathbf1}
 {\sqrt{n^{-1}\sum_j(x_j-\bar x)^2}}.
\]
There is no epsilon or affine term. A zero denominator produces $\bot$ and is
never reached. Mean, square, root and division operations use the stated grid;
all reached root arguments equal one. ReLU and argmax are exact on grid values.

\paragraph{Relation to StateBench.}
Definitions~9--10 fix the sublayer order used here. Definition~13 specifies closed
gate ranges but also inherits the gate parameterization of~\citet{gdn}.
StateBench itself uses endpoint gates in Lemmas~1b and~3, which supports an
abstract endpoint interpretation; it does not by itself resolve the inherited
activation/cast program. Its normalization inheritance includes a mean-zero
Euclidean hash construction~\citep[Appendix~C]{norm}, while the complete
fixed-grid centering/root/division path is not locally specified
in~\citet{merrill2026,li,statebench}. Theorem~\ref{thm:main} therefore states an
unconditional result in the explicit operational model above. Identifying it
with StateBench requires that model's gate and normalization interpretation to
admit the same execution. This is a specification question, distinct from
assuming or excluding finite-logit realizability in practical implementations.

\section{Rounded cyclic counters}\label{sec:counter}

\begin{lemma}[Counter]\label{lem:counter}
Let $D=Lm^2$ with $m\ge2$, $L\ge1$. Let $k$ have $m^2$ coordinates equal to $1/m$ (so $\|k\|=1$), $q=e_1$, $\alpha=\beta=1$, and let one value column hold $v=a/L$ for an integer increment $a$. If every active coordinate of the state column equals $r/(Lm)$ for an integer $r$, and the scalar intermediates are in range, then after one step every active coordinate equals $r'/(Lm)$ with
\begin{equation}\label{eq:counter}
r'=r+a+m\,\R(-r/m)=:c_m(r)+a,
\end{equation}
where $c_m(r)\in(-m/2,\,m/2]$ is the centered representative of $r\bmod m$.
\end{lemma}
\begin{proof}
In numerator units the state coordinate is $mr$. The materialized $A$ has diagonal $1-1/m^2$ and off-diagonal $-1/m^2$. A diagonal product contributes $\R\big((1-1/m^2)\,mr\big)=mr+\R(-r/m)$ since $mr$ is an integer; each of the $m^2-1$ off-diagonal products contributes $\R(-r/m)$. Summing the $m^2$ already-rounded terms gives $mr+m^2\R(-r/m)$. The write $k_iv=(1/m)(a/L)$ has numerator $ma$ exactly. Dividing by $m$ returns~\eqref{eq:counter}.
\end{proof}

Iterating~\eqref{eq:counter} from $r_0$ gives $r_t\equiv r_0+\sum_{i\le t}a_i \pmod m$ with $r_t\in(-m/2,\,m/2]+a_t$, a bounded orbit whose intermediates fit a sufficiently large fixed $p$
(Appendix~\ref{app:datatype}). This uniform-row construction differs from StateBench's two-dimensional
parity projector. Here $m^2\mid D$, so $m=2,D=25$ is not an instance.

\begin{remark}[The rescaling must move the state]
Raising $D$ to $Lm^2$ while keeping the old state $r/m$ and value $a$ breaks the orbit. At $m=3$, $D=18$, $S_i=1/3$, $a=0$: the diagonal product has numerator $\R(16\cdot6/18)=5$ and each of the eight off-diagonal products has numerator $\R(-12/18)=-1$, so the new coordinate is $(5-8)/18=-1/6$ rather than $1/3$. The certificate checks this retained-update witness separately from the
scaled identity; it does not claim a whole-network failure on a decoded word. The point matters because composing components forces a common denominator.
\end{remark}

\paragraph{Constant and input channels in the same head.} A value column receiving $mb/L$ overwrites its state to $b/L$ exactly: if the column previously held $b'/L$, the old grid numerator is $m^2b'$, so its diagonal and off-diagonal
contributions are $(m^2-1)b'$ and $-b'$, totaling zero. The write has numerator
$m^2b$, giving the new physical state $b/L$. We use this to carry the current increment $a_t/L$ and the constant $1/L$ alongside the counters, so the output projection can form arbitrary fixed rational combinations of $r_j$, $a_t$ and $1$.

\paragraph{A first-position marker from the zero state.} StateBench's Lemma~3 initializes the orbit with a beginning-of-sequence value. We avoid that token. A head with $k=\tfrac16(-5,-3,-1,-1)$ (unit norm), $q=e_1$, $\alpha=\beta=1$ and two value columns $2/D$, $1/D$ has active transition $A=I-kk^\top=\tfrac1{36}\big(\begin{smallmatrix}11&-15&-5&-5\\-15&27&-3&-3\\-5&-3&35&-1\\-5&-3&-1&35\end{smallmatrix}\big)$. Writing $T$ for its rounded numerator action, direct evaluation gives, with $u=(-2,-1,0,0)$, $v=(-3,-1,0,0)$, $w=(-1,0,0,0)$: $\R(2k)=u$, $\R(k)=w$, $T(u)=T(v)=w$, $T(w)=0$. The two columns therefore follow $0\to u\to v\to v\to\cdots$ and $0\to w\to w\to\cdots$. Reading the first coordinates $\sigma_t,\gamma_t$,
\[
B_t:=D\sigma_t-3D\gamma_t=\mathbf 1[t=1],\qquad C_t:=-D\gamma_t=1\quad(t\ge1).
\]
Both are linear in the head output. The orbit uses the tie rule at $-1/2$; under ties-away-from-zero the displayed stationary orbit is not preserved.

\begin{figure}[t]
\centering
\small
\fbox{\begin{minipage}{0.95\linewidth}
\centering
\textbf{Two complete layers, width 28; indices below are zero-based}\\[3pt]
$x_{a_t}$\ $\longrightarrow$\
$\underbrace{\LN\ \to\ \text{marker head}\ \to\ O\ \to\ \LN
\ \to\ {+x_{a_t}}\ \to\ \LN\ \to\ \mathrm{FF}+g}_{\text{layer 1: }S_0=0,\ B_t=\mathbf1[t=1]}$
\ $\longrightarrow$\ $x_{(a_t,B_t)}$\\[6pt]
$\underbrace{\LN\ \to\ \text{shifted counters}\ \to\ O\ \to\ \LN
\ \to\ {+x_{(a_t,B_t)}}\ \to\ \LN\ \to\ \mathrm{FF}+g}_{\text{layer 2: exclusive phase }c=u_{s_{t-1}}}$
\ $\longrightarrow$\ $x_{s_t}$\ $\stackrel{U}{\longrightarrow}$\ $s_t$\\[6pt]
\begin{tabular}{@{}lrl@{}}
Coordinates & Count & Layer-2 output write $y$ \\
\hline
$0$--$11$: six token-pair slots & 12 & zero \\
$12$--$17$: three phase pairs & 6 & $(c_0,-c_0,c_1,-c_1,c_2,-c_2)/4$ \\
$18$--$21$: two padding pairs & 4 & $(1,-1,1,-1)/4$ \\
$22$--$25$: unused & 4 & zero \\
$26$--$27$: shared pair & 2 & $(-1/2,+1/2)$ \\
\hline
Total & 28 & $\|y\|^2=1$, $\sum_i y_i=0$
\end{tabular}\\[4pt]
Input $x$ has one active token pair and shared $(+1/2,-1/2)$.\\
In $g=x+y$, the shared pair cancels; token, phase and padding squared norms\\
are $1/2$, $1/4$, $1/4$. The FF selector writes $x_{s_t}-g$.
\end{minipage}}
\caption{Coordinate-exact Euclidean implementation at $D=1008$, $L=112$,
$p=21$. Each FF is $W_2\relu(W_1\LN(g))$ followed by its residual addition;
its write is not normalized. In layer 1, coordinates $12$--$15$ instead hold
the two marker-label pairs, and $16$--$25$ are zero. The FF returns a canonical
label, clearing that workspace before layer 2 reuses it. Each layer has one
head. The source checker verifies these slots and the worked trace in
Table~\ref{tab:example} through the actual matrix implementation.}
\label{fig:z3}
\end{figure}



\section{A layer-norm-exact interface}\label{sec:interface}

A scalar counter readout does not by itself supply the reusable
constant-norm interface required by the following construction. The fix is to run $m$ counters whose joint state has constant norm.

\paragraph{Phase codes.} Let $u_s=\big(c_m(s+j)\big)_{j=0}^{m-1}$ for $s\in\Z_m$ and $H=\sum_j c_m(j)^2$. The $u_s$ are distinct cyclic shifts of one integer vector, with $\|u_s\|^2=H$. For $s\ne t$, Cauchy--Schwarz gives $u_s^\top u_t\le H$, and equality would force $u_s=u_t$. Thus the integer $u_s^\top u_t$ is strictly below $H$, hence at most $H-1$; equivalently $H-u_s^\top u_t=\|u_s-u_t\|^2/2\ge1$. For $m=3$: $u_0=(0,1,-1)$, $u_1=(1,-1,0)$, $u_2=(-1,0,1)$, $H=2$.

First use a marker layer and a finite recoder to retain the current label
and $B_t$. Run $m$ counter columns $j=0,\dots,m-1$ with values $v_j=(a_t+c_m(j)B_t)/L$. The marker fires once, at $t=1$, so $r_{j,1}=a_1+c_m(j)$, and thereafter each column follows~\eqref{eq:counter}. By induction,
\begin{equation}\label{eq:phase}
\big(r_{j,t}-a_t\big)_{j=0}^{m-1}=u_{s_{t-1}},\qquad s_t=\sum_{i\le t}a_i \bmod m,\ s_0=0.
\end{equation}
The head output $S_t^\top q$ thus exposes $r_{j,t}/(mL)$, $a_t/L$ and $1/L$; the output projection $O$ forms $(r_{j,t}-a_t)/(H+1)$ with integer coefficients $mL/(H+1)$ and $-L/(H+1)$ once $L$ is divisible by $H+1$.

\paragraph{Canonical carrier.} For a finite label alphabet $Z$, encode label $z$ as
\[
x_z=\tfrac12\,(e_z,\,-e_z;\ 1,\,-1),
\]
a centered unit vector: $\LN_{\mathrm E}(x_z)=x_z$. Figure~\ref{fig:z3}
and its certificate interleave each positive coordinate with its negative
partner, a fixed permutation of the displayed block notation. Its last two coordinates are a shared constant pair; $2(x_z)_{\mathrm{shared},+}=1$ supplies every constant the projections need, and $2(x_z)_{\mathrm{token},z'}=\mathbf 1[z=z']$ supplies every label indicator. Let the sublayer write be
\[
y_s=\Big(\tfrac{u_s}{H+1},\ -\tfrac{u_s}{H+1};\ \tfrac{H-1}{2(H+1)},\ -\tfrac{H-1}{2(H+1)};\ -\tfrac12,\ \tfrac12\Big),
\]
in disjoint coordinate regions (phase pair, padding pair, shared pair). Its private part has squared norm $2H/(H+1)^2+(H-1)^2/(2(H+1)^2)=\tfrac12$, so $\|y_s\|=1$, $y_s$ is centered, and $x_z^\top y_s=-\tfrac12$ through the shared pair. Hence
\[
g=x_z+y_s \quad\text{satisfies}\quad \textstyle\sum_jg_j=0,\quad \|g\|^2=1+1-1=1 .
\]
All three normalizations in the layer, $\LN(h)$, $\LN(O\,o)$ and $\LN(g)$, receive centered unit vectors and return them unchanged; in the fixed-point implementation every mean is zero and every squared Euclidean norm is exactly $1$
(the population variance is $1/n$), and the only root taken is $\sqrt{1}$. Normalization is executed, not bypassed.

\paragraph{Decoder.} The hidden unit
\[
J_{s,z}(g)=\relu\Big((H+1)\,u_s^\top g_{\mathrm{phase}}+2g_{\mathrm{token},z}-2H\textstyle\sum_{z'}g_{\mathrm{token},z'}\Big)
\]
has pre-activation $u_s^\top u_{s_{t-1}}-H+\mathbf 1[z=z_t]$, which equals $1$ when $(s,z)=(s_{t-1},z_t)$ and is $\le0$ otherwise. The $m|Z|$ units are therefore exactly one-hot without a bias. Setting the $W_2$ column for $(s,z)$ to $x_{F(s,z)}-g(s,z)$, for any finite recoding $F$, gives $h'=g+W_2J(g)=x_{F(s_{t-1},z_t)}$. Taking $F(s,z)=(z,\,s,\,s+a(z)\bmod m)$ preserves the input label and both the exclusive and inclusive prefix, which is what composition needs. The unembedding $U_r(h)=2h_{\mathrm{token},r}$ yields one logit $1$ and the rest $0$.

\paragraph{Centered unit variance.} Under $\LN_{\mathrm V}$ the fixed points are centered vectors of squared norm $n$. Uniformly scaling the Euclidean construction by $\sqrt n$ can leave no coordinate from which the tiny constants $1/D,2/D$ of the marker head are legally generated. Instead take $n=4\rho^2$, token pair $\pm1$, fixed private pairs whose squared amplitudes sum to $\rho^2-1$ and include a permanent $+1$ coordinate (a four-square decomposition supplies them), and shared pair $\pm\rho$; scale the Euclidean private write by $2\rho$ and reverse the shared pair. Use disjoint input-private and output-private coordinate regions. Then
$\|x\|^2=\|y\|^2=\|g\|^2=n$, the selector becomes $\relu\big(\tfrac{H+1}{2\rho}u_s^\top g_{\mathrm{phase}}+g_{\mathrm{token},z}-Hg_{\mathrm{const}}\big)$ with the same integer margin, and the $+1$ coordinate generates all small constants. The concrete $m=3$ network uses $n=36$, amplitudes $(2,1,1,1,1)$, shared $\pm3$.

\paragraph{Concrete $\Z_3$ network.} Two layers (marker, then phase), one head each, width $28$, feedforward width $18$, $D=1008$, $L=112$, $p=21$ for $\LN_{\mathrm E}$; width $36$, $p=24$ for $\LN_{\mathrm V}$. The variance construction needs the larger numerator range: at $p=21$,
a parameter grid numerator $3{,}048{,}192$ exceeds $2^{21}$.

\begin{table}[t]
\centering
\small
\begin{tabular}{@{}rrrlll@{}}
\toprule
$t$ & $a_t$ & $B_t$ & $(r_0,r_1,r_2)$ & $r_t-a_t\mathbf1=u_{s_{t-1}}$ & $s_t$ \\
\midrule
1 & 1 & 1 & $(1,2,0)$ & $(0,1,-1)$ & 1 \\
2 & 1 & 0 & $(2,0,1)$ & $(1,-1,0)$ & 2 \\
3 & 2 & 0 & $(1,2,3)$ & $(-1,0,1)$ & 1 \\
\bottomrule
\end{tabular}
\caption{Complete-network trace on $(1,1,2)$. The first-position bit initializes
the three offsets once. The phase is the exclusive prefix code; the decoded
residue $s_t$ is inclusive. The same scalar recurrence independently gives
$r_{j,1}=a_1+c_3(j)$ and $r_{j,t}=a_t+c_3(r_{j,t-1})$ thereafter.}
\label{tab:example}
\end{table}

\section{Composition and the exclusive prefix}\label{sec:wreath}

\paragraph{Flip-flop.} For $U_2=\{1,a,b\}$ use two state rows, $q=e_1$, key $e_0$ on the identity and $e_1$ on $a,b$, values $0,e_a,e_b$, $\alpha=\beta=1$. Row~$0$ stays zero; the identity leaves row~$1$ untouched and a reset overwrites it. With the constant channel, $O$ forms the one-hot $(1-S_a-S_b,\,S_a,\,S_b)$ including the initial identity. A separate overwrite head supplies that
constant. The recurrent update uses only $\{0,\pm1\}$; the carrier uses the rational
grid specified above. (StateBench's Lemma~1b uses $\beta=0$ on the identity; we use key selection instead so that every head in the network has $\alpha=\beta=1$.)

\paragraph{Exclusive-prefix interface.} For $W=M\wr N$ with input letters $(f_t,n_t)$ and prefix $(F_t,N_t)$, Definition~4 of~\citet{statebench} gives
\[
N_t=N_{t-1}n_t,\qquad F_t(x)=F_{t-1}(x)\,f_t(xN_{t-1}).
\]
The upper increment depends on the lower prefix \emph{excluding} the current letter. The proof of Lemma~2 proposes a feedforward lookup from represented partial
products. If these carry only the inclusive $N_t$, the following obstruction
applies to that connector.

\begin{proposition}
No function of $(f_t,n_t,N_t)$ computes $f_t(xN_{t-1})$ in general. Over $\Z_3\wr U_2$, let $f(1)=f(a)=0$, $f(b)=1$ and let $0$ denote the constant-zero function. The words $(0,a)(f,a)$ and $(0,b)(f,a)$ have the same second letter and the same inclusive prefix $N_2=a$, yet the upper increments at position~2 are $f(a)=0$ and $f(b)=1$.
\end{proposition}

The displayed pair contradicts the existence of that lookup at $x=1_N$.
It establishes a missing interface obligation for an inclusive-only proof; it does not show the lemma's depth bound is unattainable by some other construction, and we do not label the full lemma an erratum.

\paragraph{Repair: a legal one-step delay.} Run a $\Z_2$ phase component on constant increment $1$ to obtain $p_t=t\bmod 2$ as a label. A buffer head with $k_t=e_{p_t}$, $q_t=e_{1-p_t}$, $v_t=e_{N_t}$, $\alpha=\beta=1$ overwrites row $p_t$ and reads row $1-p_t$, so $S_t^\top q_t=0$ at $t=1$ and $e_{N_{t-1}}$ for $t\ge2$; using a separate overwrite companion for the current $B_t$ and constant,
adding $B_te_{1_N}$ from the marker gives $e_{N_{t-1}}$ for all $t\ge1$. All multipliers are $0$ or $1$. The decoder of \S\ref{sec:interface} then emits a canonical label carrying $(f_t,n_t,N_t,N_{t-1})$.

\paragraph{Closure.} Given complete components for $M$ and $N$ that preserve any finite payload, the $W$ component computes $N_t$, delays it, and for each fixed $x\in N$ runs a copy of the $M$ component on the increment $f_t(xN_{t-1})$, which is a finite function of the current canonical label and is realized by the same projection-plus-decoder pattern. Copies run sequentially in depth, so no two normalized carriers are ever added. The depth is finite and independent of input length; no sharp depth bound
is asserted. For a divisor $A$ of $W$ (a quotient of a submonoid $T\le W$), choose any section $\sigma:A\to T$ of the quotient map $\pi$; run the $W$ program on $\sigma(a_t)$ and decode $\pi$; multiplicativity of $\sigma$ is not needed since $\pi(\prod\sigma(a_i))=\prod a_i$. With the soluble Krohn--Rhodes decomposition into $U_2$ and cyclic-prime factors~\citep[Prop.~1, Def.~8]{statebench}, this proves Theorem~\ref{thm:main}.

\paragraph{One datatype.} Fix the program, alphabets, and a common width and head count first (padding with zero heads). Choose $D$ divisible by $36$, by every factor's $m^2$, and by every carrier, projection and normalization denominator, and let $L_m=D/m^2$ with $(H_m+1)\mid L_m$. The fixed rational ordinary products are exact after clearing their
individual denominators; $1/D$ recurrent products require the separate
primitive audit of Appendix~\ref{app:datatype}. In particular, the only intentionally rounded products are those of Lemma~\ref{lem:counter} and the marker, whose rounded maps are independent of $D$ by construction. All state sets are finite and all scalar computations have fixed length, so a finite $p$ bounds every parameter and every pre-rounding partial result for all input lengths.

\section{Exact verification}\label{sec:verify}

For each concrete instance below, a certificate builds $E,Q,K,V,O,W_1,W_2,U$, executes the state matrices with the arithmetic of \S\ref{sec:model} (including the normalization primitives, which are computed, not replaced by the identity), rejects any pre-rounding grid numerator outside $[-2^p,2^p]$, and exhausts the reachable transition system from $S_0=0$. Structural zero products are omitted while preserving the order of nonzero
products. Because the state space is finite and closed, closure certifies all
nonempty input lengths for these concrete networks. The certificate implementations were originally written from the mathematical
specification without reading the original construction implementation; this is independence of implementation, not of invention.

\begin{center}\small
\begin{tabular}{@{}llll@{}}
\toprule
Network & Convention & Width / $D$ / $p$ & States / transitions \\
\midrule
$\Z_3$ (2 layers) & $\LN_{\mathrm E}$ & 28 / 1008 / 21 & 13 / 39 \\
$\Z_3$ (2 layers) & $\LN_{\mathrm V}$ & 36 / 1008 / 24 & 13 / 39 \\
$\Z_3\wr U_2$ (7 layers, 81 letters) & $\LN_{\mathrm E}$ & 672 / 96768 / 64 & 5{,}104 / 413{,}424 \\
$\Z_3\wr U_2$ (7 layers, 81 letters) & $\LN_{\mathrm V}$ & 784 / 7056 / 64 & 5{,}104 / 413{,}424 \\
\bottomrule
\end{tabular}
\end{center}

For each $\Z_3$ arm: $390$ normalization-routine calls (including
verifier-only canonical fixed-point checks), all with exact zero mean and exact unit norm (resp.\ variance); no overflow; all $9{,}840$ nonempty words of length $\le8$ and $200$ random words of length $500$--$5000$ decoded correctly at every prefix. Short replay checks cover $73{,}812$
positions and long replay checks $538{,}724$ positions per arm, using seed
$20261108$. They replay transitions already evaluated through the full
fixed-point network; they are not extra matrix evaluations. Wreath verification
also memoizes complete local layer transitions; all $413{,}424$ global edges
are checked against a separately computed reference product. The maximum
pre-rounding grid numerators are $1{,}524{,}096$ and $9{,}144{,}576$ for the
two $\Z_3$ arms, and $14{,}046{,}068{,}736$ and $2{,}091{,}059{,}712$ for
the respective wreath arms; all are below their displayed $2^p$ bounds. The $\Z_3\wr U_2$ certificates include both reset witnesses of \S\ref{sec:wreath} and $22{,}977$ normalization calls per arm, all exact. The general-modulus algebra of \S\ref{sec:interface} (phase identities, norm identities, selector margins) was checked for $m=2,\dots,13$ ($818$ phase pairs); the proof for general $m$ is the induction above, not the enumeration.

\paragraph{Which assumptions carry the weight.}
Re-running the $\Z_3$ certificates with ties-to-even, or with a floored square
root, passes (every root argument is $1$). Fused-exact evaluation instead
computes the rank-one affine recurrence with exact intermediates and one cast
per state output. It first breaks the marker canonical invariant on $[0,0]$,
before the final cyclic readout. Ties-away-from-zero breaks the marker orbit.
These two negative witnesses concern intermediate invariants, not demonstrated
incorrect final modulo answers. The passing ties-even arm shows that the
baseline tie rule is not necessary for this particular network.

\paragraph{Complete-output rank-one sensitivity.}
A separate alternative evaluates the rank-one update with individually rounded
products. Write $K_i,S_{ij},V_j$ for integer grid numerators and compute
\begin{align*}
 d_j&=\sum_{i=0}^{q-1}\R(K_iS_{ij}/D),\\
 S'_{ij}&=\bigl(S_{ij}-\R(K_id_j/D)\bigr)+\R(K_iV_j/D).
\end{align*}
Here $q$ is the number of active key rows. Sums are checked left to right;
subtraction precedes write addition. The matrix $I-kk^\top$ is not materialized.
All other operations, normalization checks and overflow checks stay unchanged.
This is not the fused-exact evaluator above. Complete runs from zero memory
under both normalization constructions give the same table:
\begin{center}
\begin{tabular}{@{}lll@{}}
\toprule
Input word & Materialized outputs & Rounded rank-one outputs \\
\midrule
$[1,1]$ & $[1,2]$ & $[1,1]$ \\
$[1,2]$ & $[1,0]$ & $[1,2]$ \\
$[1,0]$ & $[1,1]$ & $[1,0]$ \\
\bottomrule
\end{tabular}
\end{center}
The certificate executes twelve complete network runs, with no overflow.
The mechanism is explicit: for a uniform counter column of numerator $mr$,
each rounded product in $k^\top S$ has numerator $r$, giving $d=m^2r$;
the correction has numerator $mr$ and cancels the retained state exactly.
For the marker, $\sum_i\R(k_i u_i)=3$ and $\R(3k)=u$, while
$\sum_i\R(k_i w_i)=1$ and $\R(k)=w$. Thus the alternative repeats its
first-position column states, reports $B_t=1$, and reinitializes the phase
at each step. This explains the current-increment outputs in the table.
These results concern the specified evaluators, not production kernels.

\section{Discussion}

\paragraph{Relation to prior results.}
\citet[Theorem~1; Lemmas~1b--3]{statebench} supply aperiodic and
rounded-parity baselines for GDN$^+$; their Theorem~2 gives a soluble lower bound
for log-precision averaging hard attention. We construct arbitrary-modulus
counters, normalized payload-preserving interfaces and complete soluble closure
in the explicit materialized model. Our exclusive-prefix obstruction concerns
a particular inclusive-only connector, not the attainability of the cited
composition lemma's depth bound.
\citet[Theorem~1; Appendix~A.4]{grazzi} prove a nonnegative-spectrum parity
obstruction under a different exact-intermediate/casting convention. Separate
scalar casts change the transition, so there is no contradiction under shared
assumptions. \citet[Theorems~1--2; Appendix~B.1]{siems} obtain broad group and
all-regular expressivity with DeltaProduct, including one update per head at
sufficient depth, with coefficients up to $2$, abstract parameter maps and
decoder, and task-dependent initialization or a BOS token. These results do not
supply the present zero-state, constant-gate, individually rounded interfaces.
\citet[Appendix~J]{nowak} provide the arithmetic-sensitive algebraic framework
and diagonal-recurrence cases. Unsigned-integer wrap-around differs from the
present signed fixed-grid rounding. \citet[Lemmas~18, 20]{sarrof} already
implement previous-token access and delayed composition for SSMs. Neither
rounding-induced periodicity, Krohn--Rhodes decomposition nor delayed composition
is claimed as new here.

NFSM \citep[Definitions~2.1--2.2]{nfsm} requires correct decoding along every
trajectory with arbitrary per-step perturbations bounded by a positive radius.
Our theorem concerns one deterministic rounding program, not that robustness
guarantee. Its rounded state updates are nonlinear maps, so the affine
robustness obstruction does not contradict this result. Complex KDA
\citep[Theorems~3, 5; Appendix~B.4]{ckda} uses signed channel gates and extended
delta coefficients; its constructive arithmetic uses exact datatypes, potentially
containing algebraic irrationals. Neither those results nor ours transfer
automatically to floating-point execution.

\paragraph{Construction cost.}
The cyclic head represents $m^2(m+2)=\Theta(m^3)$ active recurrent entries.
Naive dense evaluation of its retained product uses
$m^4(m+2)=\Theta(m^5)$ scalar multiplications per token:
$m^2(m+2)$ dot products of length $m^2$. These are active-block counts, not
minimal storage, bit complexity or full padded-network cost. Composition adds
$|N|$ sequential upper blocks, a clock, delay and finite controllers to the
lower block; no optimal-depth or polynomial-size bound is asserted.

\paragraph{Limitations.}
Exact StateBench identification remains open. This is an existence bound without
an efficient-size guarantee, upper bound, empty-word result, or consequence for
nonzero epsilon, practical kernels or trained models. Finite certificates
validate concrete networks; the general theorem is algebraic. The inclusive-prefix
witness does not refute Lemma~2's statement or depth bound.

\paragraph{Reproducibility.}
Repository: \url{https://github.com/TarangKhanna/gdn-plus-soluble}.
The verification implementation is fixed at commit
\href{https://github.com/TarangKhanna/gdn-plus-soluble/tree/b4c080f9115598af59d0e0543ee7fbb704283e08}{\nolinkurl{b4c080f9115598af59d0e0543ee7fbb704283e08}}.
The accompanying LaTeX source package contains its exact export at
\texttt{anc/gdn\_plus\_verification\_code.zip}; the repository's \texttt{paper/}
directory contains this revision's PDF and source. The
\href{https://doi.org/10.5281/zenodo.24057440}{original Zenodo preprint}
(DOI: 10.5281/zenodo.24057440) and the
\href{https://github.com/TarangKhanna/gdn-plus-soluble/tree/871536d6037326b5ddf222763a4c58d6e7e3f646}{immutable original code baseline}
retain the earlier provenance. The current implementation exposes the
constant-gate wreath construction directly and adds the complete-output
rank-one certificate; original expected results remain unchanged.
From a clean extraction of the code archive, run
\texttt{python3 -B certs/run\_all.py}, without \texttt{-O}.
The full suite passed on macOS with Python 3.14.4; no new Linux or Python 3.10
verification is claimed. Code is standard-library only. The general theorem
rests on the proof appendices, not the finite certificates. AI assistance in
construction, verification, review and drafting is disclosed on the first page;
the author retains responsibility for the claims and final revision.

\begin{thebibliography}{11}\small
\setlength{\itemsep}{0pt}
\paperbibitem{Grazzi et al.(2025)}{grazzi} R.~Grazzi, J.~Siems, A.~Zela, J.~Franke, F.~Hutter, M.~Pontil. Unlocking state-tracking in linear RNNs through negative eigenvalues. ICLR 2025; arXiv:2411.12537v4.
\paperbibitem{Li et al.(2026)}{li} Y.~Li, A.~Svete, A.~Sabharwal, W.~Merrill. Efficiently representing algorithms with chain-of-thought transformers. arXiv:2606.19697, 2026.
\paperbibitem{Merrill et al.(2026)}{merrill2026} W.~Merrill, H.~Jiang, Y.~Li, A.~Lin, A.~Sabharwal. Why are linear RNNs more parallelizable? arXiv:2603.03612, 2026.
\paperbibitem{Merrill and Sabharwal(2023)}{norm} W.~Merrill, A.~Sabharwal. The expressive power of transformers with chain of thought. arXiv:2310.07923v5.
\paperbibitem{Nowak et al.(2026)}{nowak} F.~Nowak, R.~Cotterell, R.~Boumasmoud. An algebraic view of the expressivity of recurrent language models. ICML 2026; arXiv:2606.01765v3.
\paperbibitem{Petty et al.(2026)}{statebench} J.~Petty, W.~Merrill, A.~Sabharwal. A fine-grained hierarchy and benchmark for state tracking in language models. arXiv:2611.03094v1, 2026.
\paperbibitem{Sarrof et al.(2024)}{sarrof} Y.~Sarrof, Y.~Veitsman, M.~Hahn. The expressive capacity of state space models: a formal language perspective. NeurIPS 2024.
\paperbibitem{Siems et al.(2025)}{siems} J.~Siems, T.~Carstensen, A.~Zela, F.~Hutter, M.~Pontil, R.~Grazzi. DeltaProduct: improving state-tracking in linear RNNs via Householder products. NeurIPS 2025; arXiv:2502.10297v7.
\paperbibitem{Yang et al.(2024)}{gdn} S.~Yang, J.~Kautz, A.~Hatamizadeh. Gated delta networks: improving Mamba2 with delta rule. arXiv:2412.06464, 2024.
\paperbibitem{Brandoit et al.(2026)}{nfsm} J.~Brandoit, A.~Fyon, T.~Braipson, T.~Clara, F.~De~Geeter, P.~Sacr\'e, D.~Ernst, G.~Drion. Length-independent state tracking under a parallel scan. \href{https://arxiv.org/abs/2609.32447v1}{arXiv:2609.32447v1}, 2026.
\paperbibitem{Siems et al.(2026)}{ckda} J.~Siems, R.~Grazzi, K.~P\"oppel, J.~Singh, A.~Zela, T.~Carstensen, J.~Jitsev, F.~Hutter, V.~Cevher, A.~Orvieto, A.~Klein. Complex KDA: understanding and enhancing the expressivity of Kimi Delta Attention. \href{https://arxiv.org/abs/2609.24797v1}{arXiv:2609.24797v1}, 2026.
\end{thebibliography}

\appendix
\section{Primitive constructions and exact arithmetic}\label{app:primitives}
This appendix gives the local implementation steps used by the induction.
A label always denotes an entire finite tuple, encoded with one active token
pair. Enlarging the tuple alphabet enlarges width; it does not concatenate
several independently normalized labels.

\subsection{Projection and finite recoding}
For Euclidean canonical input $x_z$, a linear map with column $2b_z$ at the
positive token coordinate returns any prescribed rational vector $b_z$ on
label $z$. The shared positive coordinate supplies a constant by weight twice
that constant. After choosing the grid as below, each individual product is
representable; no affine bias is used.

A finite function $F:Z\to B$ can be computed as a complete layer. Use the
unit-key overwrite head $k=q=e_0$, $v=(e_{F(z)},1)$, so the readout is the
current one-hot and constant, regardless of previous memory. Let its output
projection emit a private pair $(e_b/2,-e_b/2)$ and shared pair $(-1/2,1/2)$.
The residual contains the input-token pair and output-private pair with the
shared pair canceled. For each $(z,b)$, the bias-free selector
\[
 J_{z,b}(g)=\relu\!\left(2g_{{\rm token},z}
       +2g_{{\rm output},b}-2\sum_{z'}g_{{\rm token},z'}\right)
\]
is exactly one for the actual pair and zero otherwise. Its $W_2$ column is
$x_{G(z,b)}-g(z,b)$, where $G$ can retain any original payload together with
$F(z)$. This is a legal finite controller, not an external runtime table.

The marker uses the same output-pair geometry with $b=B_t$. Its two value
columns are generated from the shared $+1/2$ by weights $4/D$ and $2/D$.
For unit-variance inputs the permanent $+1$ generates them directly.
The marker supplies $B_t$ at every position; a recoder can retain it with the
entire incoming label. Subsequent layers can recompute or preserve it.

\subsection{Cyclic initialization and boundedness}
Let the finite label determine $a_t\in\{0,\ldots,m-1\}$. Define
$c_m(r)=r+mR(-r/m)$. This is periodic modulo $m$, congruent to $r$, and takes
values in $(-m/2,m/2]$. For the shifted column initialized with the marker,
\[
 r_{j,1}=a_1+c_m(j),\qquad
 r_{j,t}=a_t+c_m(r_{j,t-1})\quad(t\ge2).
\]
Assume $r_{j,t-1}=a_{t-1}+c_m(s_{t-2}+j)$. Periodicity and congruence give
$c_m(r_{j,t-1})=c_m(s_{t-1}+j)$, proving the claimed exclusive phase. This
also proves that the raw integer states range over a fixed finite set,
independent of length. At the first step the same identity holds with $s_0=0$.
The companion columns with values $ma_t/L$ and $m/L$ supply the current input
and constant by exact overwrite, as calculated in the main text.

For completeness, the marker's row action is
$T(z)_i=\sum_jR(A_{ij}z_j)$ in integer numerator units. Direct substitution
in its displayed $4\times4$ matrix gives
\[
 R(2k)=u,\quad R(k)=w,\quad T(u)=w,\quad T(v)=w,\quad T(w)=0,
 \qquad v=u+w.
\]
Thus its memory invariant is exactly the two finite column orbits given in the
main text; the initialization correction is available at all lengths.

\subsection{The concrete modulo-three Euclidean carrier}
The generic phase carrier uses denominator $H+1$ and is convenient for the
proof. The width-28 executable certificate uses a different, equally exact
private geometry: $c=u_s\in\{(0,1,-1),(1,-1,0),(-1,0,1)\}$ and
\[
 y=\big(c/4,-c/4;1/4,-1/4;1/4,-1/4;-1/2,1/2\big).
\]
The first two padding pairs have total squared norm $1/4$, the phase pairs
have squared norm $1/4$, and the shared pair has squared norm $1/2$.
After adding the input token pair, the shared pair cancels and the residual
again has squared norm one. The head reads $r_j/(3L),a/L,1/L$; output weights
$3L/4,-L/4$ form $(r_j-a)/4$. Padding and shared coordinates use weights
$\pm L/4,\pm L/2$. For label $\ell=(a,B)$ and padding coordinate $1/4$,
\[
 J_{s,\ell}(g)=\relu\!\left(4u_s^\top g_{\rm phase}
             +2g_{{\rm token},\ell}-8g_{\rm pad}\right)
\]
is exactly the joint indicator. Its $W_2$ column is
$x_{(s+a)\bmod3}-g_{s,\ell}$, where $g_{s,\ell}$ is the reached residual for
that joint label. The residual addition therefore returns
$g_{s,\ell}+(x_{(s+a)\bmod3}-g_{s,\ell})=x_{(s+a)\bmod3}$,
clearing the phase and padding coordinates.
All 18 hidden units are explicit and bias-free. The last unembedding is
$U_r(h)=2h_{{\rm token},r}$, giving a unique unit logit.
The variance certificate has its own carrier and selector constants; it is
not obtained merely by changing the normalizer in this Euclidean network.

\section{Variance-one geometry at arbitrary width}\label{app:variance}
Fix the maximum number $K$ of joint labels and the maximum number $J$ of
private-output coordinate slots in the finite program. Choose an integer
$r\ge2$ so that $n=4r^2$ accommodates all token, fixed-private, output-private
and shared slots, with these regions disjoint. By the four-square theorem,
write $r^2-2=b_1^2+\cdots+b_4^2$. The canonical input has one active token pair
$\pm1$ among $K$ pairs, fixed-private pairs of amplitudes
$1,b_1,\ldots,b_4$, and shared pair $\pm r$. The first fixed pair includes a
permanent $+1$. Its input-private squared norm is
\[
 2+2(1+b_1^2+\cdots+b_4^2)=2r^2,
\]
and its shared squared norm is $2r^2$, giving $\|x_z\|^2=n$ and zero mean.
Only finitely many slots are needed, while $4r^2$ can be arbitrarily large.

Scale any Euclidean private write (squared norm $1/2$) by $2r$, and use the
opposite shared pair $\mp r$. Then $\|y\|^2=n$ and $x_z^\top y=-2r^2$.
The residual therefore has squared norm $n$, population variance one and
zero mean. Its permanent $+1$ input coordinate survives. The normalization
root argument is one, even though the squared Euclidean norm is $n$.

For a paired one-hot private output of amplitude $r$, the selector is
\[
 \relu(g_{{\rm token},z}+g_{{\rm output},b}/r-g_{{\rm const}}).
\]
For a cyclic private output, it is
\[
 \relu\!\left(\frac{H+1}{2r}u_s^\top g_{\rm phase}
                    +g_{{\rm token},z}-H g_{{\rm const}}\right).
\]
Both reduce to the integer-margin selectors in the Euclidean construction.
Weights $W_2=x_{\rm target}-g$ use the actual variance-one vectors. QKV
projections use token amplitude one; constants, however small, can use the
permanent $+1$. All parameters are rational after fixing integer $r$.

\section{Payload preservation, delay and algebraic closure}\label{app:closure}
\subsection{Payload-preserving components}
A component accepts a joint input label $(z,p)$, where $p$ is any finite
payload. Lift its QKV projections to depend only on the old label $z$.
The recurrent trajectory is unchanged. In the final conjunction selectors,
index by the entire actual input label $(z,p)$ and the exposed finite state.
Choose $W_2$ to return the desired output together with $(z,p)$ in one
canonical label. This works for every primitive above. It also lifts each
layer of a composed block: widths are chosen for the finite enlarged alphabets
before choosing the datatype. No factorized-amplitude payload interface is
assumed.

\subsection{Flip-flop and one-step delay}
For $U_2$ use its multiplication with $a,b$ as right-zero resets. Keys select
an inert row on identity and the active row on either reset; all gates remain
one. The active row is zero before any reset, then holds one of $e_a,e_b$.
A separate overwrite companion head, $k=q=e_0,v=1$, supplies the current
constant. The output projection forms $(1-S_a-S_b,S_a,S_b)$, including the
identity state, and the finite recoder returns the inclusive product while
retaining the full input payload. Companion heads are counted in the common
finite head count, rather than assumed to be available from the reset head.

After any lower block, run a cyclic-two block on constant increment one,
preserving the current lower product and input, to expose $p_t=t\bmod2$.
A fresh buffer head uses $k_t=e_{p_t},q_t=e_{1-p_t},v_t=e_{N_t}$. Its two
rows are initially zero. At time $t$ the selected row is overwritten and the
opposite row, most recently written at $t-1$, is read. Hence the readout is
zero for $t=1$ and $e_{N_{t-1}}$ afterwards, by induction. A separate
current-value overwrite companion $k=q=e_0,v=(1,B_t)$ supplies the constant
and first-position correction. Form $e_{N_{t-1}}$ by adding $B_te_{1_N}$.
The same paired carrier and recoder preserve all current and delayed labels.
These companion channels do not rely on the buffer retaining current values.

\subsection{Wreath products and divisors}
Suppose blocks for $M$ and $N$ compute their full-alphabet prefix products and
preserve every finite payload. Compute $N_t$, then $N_{t-1}$ as above, retaining
the current $(f_t,n_t)$. For each fixed $x\in N$, a finite controller computes
$f_t(xN_{t-1})$ as an input label to a fresh $M$ block. Run these $|N|$ blocks
sequentially in depth, preserving the lower states and previously computed
upper products. Their fresh zero memories do not interfere. By induction on
positions, the $x$th result is
\[
 F_t(x)=\prod_{i=1}^t f_i(xN_{i-1}),
\]
which, with $N_t$, is precisely the prefix product in $M\wr N$.
All additional components have fixed finite depth; no optimal sum-depth bound
is required for the theorem.

If $A$ is a homomorphic image of a submonoid $T\le W$, choose a set-theoretic
section $\sigma:A\to T$ of the surjective homomorphism $\pi:T\to A$. Run the
full-alphabet $W$ block on $\sigma(a_t)$. Every reached product lies in $T$,
and
\[
 \pi\!\left(\prod_{i=1}^t\sigma(a_i)\right)=\prod_{i=1}^t a_i.
\]
Only $\pi$ is multiplicative; the section need not be. Implement the section
and quotient decoding with the finite recoders above, extending the decoder
arbitrarily outside $T$. The soluble Krohn--Rhodes decomposition into cyclic
prime groups and $U_2$ factors, as stated in StateBench Proposition~1
(see also the classical decomposition in \citep[Appendix~I, Theorem~I.3]{nowak}), now
completes a structural induction. The trivial factor is a constant recoder.

\section{A common datatype and all-length bounds}\label{app:datatype}
First fix the finite algebraic program, every joint alphabet and coordinate
order, and a sufficiently large common width and head count. Unit carriers
pad with zero coordinates. Variance carriers use Appendix~\ref{app:variance}.
Choose the common width divisible by the head count and large enough
that every head block fits its active rows and value columns (for variance
width $4r^2$, choose $r$ a sufficiently large multiple of that count).
Pad unused heads with unit key and zero query/value, so they have zero output
and memory; do not normalize a dummy head in isolation. Fix feedforward width
large enough for all finite conjunction selectors.

All carriers, finite-controller targets, selector margins and normalization
intermediates now have fixed rational values, independent of $D$. Choose $D$
to clear every denominator of these values and their individual ordinary
products, including coordinates divided by $n$, coordinate squares and squares
divided by $n$. Include factors $36$ and $m^2$ for each cyclic modulus, and
require $L_m=D/m^2$ divisible by the fixed denominator factors of its output
coefficients. The finite list has a common multiple. Only after fixing this
$D$ do we set the $L_m$ and instantiate the grid-dependent parameters.
The theorem quantifies existentially over $D,p$ and those parameters; it does
not hold them fixed while arbitrarily changing the grid. The following audit
separates intentionally rounded operations from ordinary exact ones.
\begin{enumerate}
\setlength{\itemsep}{2pt}
\item Canonical projections, finite selectors, FF writes, binary reset/buffer
states and the norm calculations have fixed rational operands/products.
Clearing their individual denominators makes them exact, including both mean
and variance evaluation orders checked by the certificates.
\item Cyclic keys and materialized transition entries are exact once $m^2\mid D$.
State numerators are $mr$; retained products have the intentionally rounded
numerators derived in Lemma~\ref{lem:counter}, independent of $D$.
Their tiny writes have integer numerator $ma$. Overwrite companions cancel
exact integer retained terms.
\item Marker transition entries have denominator $36$, and state numerators
belong to $\{0,u,v,w\}$. Its products with these $1/D$ states deliberately
round, with the same numerator map at every allowed $D$. Values $1/D,2/D$
are generated exactly from the canonical constants; clearing denominators is
not asserted to remove the marker's intentional rounding.
\item Output coefficients proportional to $D$ or $L_m$ multiply state values
proportional to $1/D$, producing the already fixed rational carriers.
Their parameter grid numerators and individual product denominators impose
only finitely many additional divisibility conditions. There is no requirement
that a denominator depending recursively on $D$ divide $D$.
For example, a coefficient $\gamma D$ multiplying a state $z/D$ yields the
fixed rational $\gamma z$. Its parameter numerator is $\gamma D^2$;
clearing the fixed denominator of $\gamma$ suffices. Tiny marker values are
generated from the canonical $+1/2$ (Euclidean) or permanent $+1$ (variance)
coordinate, as in Appendix~\ref{app:primitives}; their intentional rounding
is retained rather than eliminated.
\end{enumerate}
\paragraph{Coverage under composition.}
Payload lifting replicates canonical projection columns and enlarges finite
selector indices; it introduces only the products of case~1. Finite recoding,
including sections, quotient decoding and the trivial factor, uses the overwrite
head and residual-corrected FF of Appendix~\ref{app:primitives}: case~1.
A parity clock reuses cases~2 and~4 and its marker uses case~3. Reset and delay
heads use binary keys, queries, values and memories; their companion channels,
first-position correction and carrier projections use cases~1 and~4.
Wreath composition sequentially connects canonical labels and makes finitely
many copies of these blocks; it does not multiply their memories together.
Common-width padding adds zeros and inert unit-key heads, covered by case~1;
variance padding uses the fixed rational geometry of Appendix~\ref{app:variance}.
In every block, gate factors are one, centering and squared-statistic products
are in case~1, the root argument is exactly one, and division is by one. Residual
and dot-product sums are ordered finite additions; ReLU, coordinate selection
and the final argmax introduce no new rounded product. This structural induction
covers the enlarged finite alphabets before $D$ is chosen.

Each primitive has a finite
memory invariant and finite canonical input alphabet. Their sequential product
therefore has a finite reachable set, independent of input length. At fixed
$D$ there are finitely many parameters and finitely many exact pre-round scalar
results and ordered partial sums in one transition. Let $B$ bound their absolute
physical values. Choose finite $p$ with $2^p>DB$. Every transition stays inside
the datatype range for every input length. This is an existence bound, not a
precision-optimality or polynomial-size claim. The bound is chosen after $D$
and includes parameters and ordered partial sums, not just final state entries.
It provides no perturbation margin or floating-point equivalence.

\end{document}
